\section{A Bias--Freedom Model}
\label{sec:theory}

% 证据：实验 6a（theory/model.py、sweep.py、meanfield.py）；T11 报告（bias_multi.py，检验 B、D）
% 不能写："理论预言了 K0.5"；区间比例 1/q 在 FJSP 上未定量验证（只在第 IX 节写定性结论）

The density $\dens$ trades bias for freedom.
This section explains the trade-off with a model that is small enough to analyze and to simulate exhaustively.
The model answers two questions: how many slots are needed, and how the best number changes with the evaluation budget.

\subsection{Rule-Distance Model}
\label{sec:theory:model}

A solution is a string $\mathbf{z} \in \{0, \dots, \nval-1\}^{\npos}$ of $\npos$ decisions.
Value $0$ is the rule's choice, so the rule's solution is the all-zero string.
The target $\mathbf{t}$ differs from the rule in $\dist$ positions, where it takes nonzero values.
The objective is the Hamming distance to the target, $f(\mathbf{z}) = |\{j : z_j \ne t_j\}|$, to be minimized.
Here $\dist$ is the \emph{rule distance}: the number of decisions in which a good solution departs from the rule.
The model ignores cascades, \ie a changed decision does not change the meaning of later ones.
Section~\ref{sec:repr} showed that cascades do not explain performance in our decoders.

We compare two encodings with genes in $[0,1)$.
\begin{itemize}
  \item \textbf{Per-step (PS)}, the abstraction of \dec{0}--\dec{5}: one gene per decision, $z_j = \lfloor g_j \nval \rfloor$.
  \item \textbf{Sparse deviation (SP$_\K$)}, the abstraction of \SDD: $\K$ slots $(p_k, a_k)$ as in Section~\ref{sec:method:genotype}.
  Slot $k$ is active if $a_k \ge 0.6$, so it is active with probability $\qon = 0.4$ under a random genotype.
  It sets position $\lfloor p_k \npos \rfloor$ to the nonzero value $1 + \lfloor u_k (\nval-1) \rfloor$ with $u_k = (a_k-0.6)/0.4$.
  Later slots overwrite earlier ones.
\end{itemize}

\subsection{Three Properties of Sparse Deviation}
\label{sec:theory:props}

The first property describes where random genotypes start.

\begin{proposition}[Initial distance]
\label{prop:init}
For a uniform random genotype, let $\pi = 1 - (1 - \qon/\npos)^\K$ be the probability that a given position is deviated. Then
\begin{equation}
  \mathbb{E}[f_{\mathrm{SP}}] = (\npos - \dist)\,\pi + \dist \left(1 - \frac{\pi}{\nval-1}\right).
\end{equation}
For PS, $\mathbb{E}[f_{\mathrm{PS}}] = \npos (1 - 1/\nval)$.
\end{proposition}

\begin{IEEEproof}
Each slot hits a given position with probability $\qon/\npos$, independently of the other slots, so the position is deviated with probability $\pi$.
The last write at a deviated position is uniform over the $\nval-1$ nonzero values.
A non-target position is wrong exactly when it is deviated.
A target position is right exactly when it is deviated and receives the right value, which has probability $\pi/(\nval-1)$.
\end{IEEEproof}

For small $\K$, $\pi \approx \qon\K/\npos$ and $\mathbb{E}[f_{\mathrm{SP}}] \approx \dist + \qon\K$.
Random genotypes start near the rule, about $\qon\K$ wrong deviations away from it.
Under PS they start at distance $\npos(1-1/\nval)$, whatever the rule is; for $\npos = 100$ and $\nval = 4$ this is 75.

The second property limits what a given $\K$ can reach.

\begin{proposition}[Freedom floor]
\label{prop:floor}
Every genotype of SP$_\K$ deviates at no more than $\K$ positions, so $\min_{\genome} f_{\mathrm{SP}} \ge \max(0, \dist - \K)$.
\end{proposition}

If $\K < \dist$, no amount of search reaches the target.
The search stalls at least $\dist - \K$ decisions away from it.

The third property counts how many genotypes represent the target.

\begin{proposition}[Redundancy of the target]
\label{prop:red}
The probability that a uniform random genotype of SP$_\K$ decodes exactly to $\mathbf{t}$ is
\begin{equation}
  r(\K) = \frac{\dist!}{(\nval-1)^{\dist}} \sum_{J=\dist}^{\K} \binom{\K}{J} \qon^J (1-\qon)^{\K-J} \npos^{-J}\, \mathcal{S}(J, \dist),
\end{equation}
where $\mathcal{S}(J,\dist)$ is the Stirling number of the second kind. In particular, $r(\K) = 0$ for $\K < \dist$.
\end{proposition}

\begin{IEEEproof}
Let $J$ be the number of active slots; $J$ is binomial with parameters $\K$ and $\qon$.
The decoded string equals $\mathbf{t}$ exactly when the $J$ active slots fall only on the $\dist$ target positions,
cover every one of them, and the last write at each one has the right value.
Of the $\npos^J$ equally likely placements, $\dist!\,\mathcal{S}(J,\dist)$ map the slots onto the $\dist$ positions surjectively.
Each of the $\dist$ last writes is right with probability $1/(\nval-1)$.
\end{IEEEproof}

\begin{corollary}
\label{cor:kstar}
Keeping only the leading term $J = \dist$, $r(\K+1)/r(\K) = (\K+1)(1-\qon)/(\K+1-\dist)$, which is at least one exactly when $\K + 1 \le \dist/\qon$.
The redundancy of the target is therefore largest near $\K^\ast \approx \dist/\qon$.
\end{corollary}

The exact maximizers agree with the corollary.
For $\npos = 100$ and $\nval = 4$, they are $\K^\ast = 12$, 26 and 55 for $\dist = 5$, 10 and 20, against $\dist/\qon = 12.5$, 25 and 50.
For PS the target has redundancy $\nval^{-\npos}$, independent of the rule.

Together, the three properties describe the trade-off.
With $\K < \dist$ the target is out of reach.
With $\K$ near $\dist/\qon$ the target is represented by the most genotypes.
Beyond that, random genotypes deviate at too many wrong positions, and both the start and the redundancy get worse.

\subsection{Behavior Under a Fixed Budget}
\label{sec:theory:budget}

The properties above describe the search space, not the search.
To see which $\K$ is best for a given number of evaluations, we run a GA with the operators of our NSGA-II:
SBX ($\eta = 20$, probability 0.9), polynomial mutation ($\eta = 20$, rate $1/n$), population 30 and truncation selection.
We use $\npos = 100$, $\nval = 4$, nine values of $\dist$ between 0 and 80, eleven values of $\K$ between 5 and 100, and 20 runs per setting.

Table~\ref{tab:theory-budget} shows two regimes.
\begin{itemize}
  \item \textbf{Small budget (465 evaluations).} The best $\K$ is close to $\dist$ or below it.
  \item \textbf{Large budget (10\,000 evaluations).} The best $\K$ moves up to about $\dist/\qon$, where the target is most redundant.
\end{itemize}
A fixed $\K = 25$ is good only in between: with $\dist = 40$ and 10\,000 evaluations it ends at distance 29.7, while $\K = 100$ reaches 12.1.
PS is far behind in all settings: its final distance is above 40 even after 10\,000 evaluations.

\begin{table}[t]
  \centering
  \caption{Best number of slots $\K$ in the rule-distance model ($\npos = 100$, $\nval = 4$, $\qon = 0.4$). Entries: best $\K$ (final distance). Mean of 20 runs.}
  \label{tab:theory-budget}
  \setlength{\tabcolsep}{4pt}
  \begin{tabular}{@{}rrccc@{}}
    \toprule
    $\dist$ & $\dist/\qon$ & 465 evals & 2000 evals & 10\,000 evals \\
    \midrule
    5  & 12.5 & 5 (4.0)       & 15 (3.5)  & 20 (2.4) \\
    10 & 25   & 5, 10 (8.6)   & 25 (6.8)  & 50 (3.8) \\
    20 & 50   & 10 (17.2)     & 30 (13.6) & 75 (6.8) \\
    40 & 100  & 20, 30 (35.1) & 60 (27.1) & 100$^\dagger$ (12.1) \\
    \bottomrule
    \multicolumn{5}{@{}l}{\footnotesize $^\dagger$Largest value tested.}
  \end{tabular}
\end{table}

\subsection{Why the Best Number of Slots Grows With the Budget}
\label{sec:theory:meanfield}

A mean-field model makes the mechanism visible.
We split the $\K$ slots into $w$ wrong active slots, $c$ correct active slots, and $i = \K - w - c$ inactive slots, so that $f = w + (\dist - c)$.
An offspring changes about one gene, so a given slot is touched with probability about $1/\K$.
With elitist selection, only changes that lower $f$ survive, and there are three of them:
\begin{align}
  \text{clean:}  &\quad w \to w-1        &&\text{rate } \alpha\, \tfrac{w}{2\K}, \\
  \text{repair:} &\quad c \to c+1        &&\text{rate } \alpha\, \tfrac{i}{2\K} \cdot \tfrac{\dist - c}{\npos(\nval-1)}, \\
  \text{move:}   &\quad w-1,\ c+1        &&\text{rate } \beta\, \tfrac{w}{2\K} \cdot \tfrac{\dist - c}{\npos(\nval-1)}.
\end{align}
The initial values of $w$ and $c$ follow from Proposition~\ref{prop:init}.
The constants $\alpha$ and $\beta$ depend only on the operators.
Fitted on the runs with 465 and 2000 evaluations ($\alpha = 0.05$, $\beta = 1.5$),
the model predicts the final distances after 10\,000 evaluations with a correlation of 0.994.

The two regimes follow from the two kinds of events.
Every wrong active slot costs one unit of distance and must be cleaned.
There are about $\qon\K$ of them at the start, so their cost grows with $\K$ and is paid first.
A repair needs an inactive slot that lands on a missing position with the right value.
When $\K$ is close to $\dist$, the inactive slots run out as the search approaches the target, and repairs slow down.
A small budget is spent mostly on cleaning, which favors few slots.
A large budget is limited by repairs, which favors spare slots.

\subsection{Several Good Solutions}
\label{sec:theory:multi}

A Pareto front contains many good solutions, and they lie at different distances from the rule.
We extend the model to $t$ targets at distances $\dist_1, \dots, \dist_t$ and run NSGA-II on the vector of distances to the targets.
Performance is the mean, over targets, of the smallest distance reached.

The best $\K$ follows the mean distance rather than the smallest one.
For example, with targets at distances 2 and 40 and 2000 evaluations, the best $\K$ is 30;
for the near target alone it would be 5.
Over all target sets, the best $\K$ is up to 7.5 times the smallest distance at 465 evaluations and up to 15 times at 2000 evaluations.
A density chosen to reach the nearest good solution is therefore too small for a front.

One more factor raises the best $\K$.
In the base model an unnecessary deviation costs as much as a wrong decision.
If it costs less, the best $\K$ grows: when the cost falls to 0.1, doubling $\K$ from 25 to 50 no longer hurts at $\dist = 5$.
A redundant decoder lowers this cost, since many deviations do not change the solution.

\subsection{Consequences for the Method}
\label{sec:theory:consequences}

The model gives three rules of thumb, which the experiments test.
\begin{enumerate}
  \item The density must be at least as large as the typical rule distance; otherwise the search stalls (Proposition~\ref{prop:floor}).
  \item Larger budgets call for more slots (Table~\ref{tab:theory-budget}).
  \item In multi-objective search the relevant distance is the mean distance to the front, not the distance to the nearest good solution (Section~\ref{sec:theory:multi}).
\end{enumerate}
The model does not give the best density for a real instance, because $\dist$ is unknown before search.
This is why \SDD chooses $\dens$ on training instances (Section~\ref{sec:method:k}).
